\documentclass{ltjsbook}
\usepackage{docmute} 
\usepackage{../../myStyle} 

\begin{document}
\chapter{データ空間}
\label{L06_DataSpace}
ベクトルを空間的なイメージとして理解できたところで，これをデータ解析に応用していきましょう。
私たちが考えたいのはデータのベクトルであり，データの空間だったわけです。これに具体的なモデルを入れて考えていきましょう。
\section{線形回帰を行列と空間で理解する}
セクション\ref{LinearModelForm}(Pp.\pageref{LinearModelForm})でみたように，回帰分析はデータベクトルと係数ベクトルを使って，
\[ \bm{y} = \bm{x\beta} + \bm{e} \]
と表せるのでした。具体例として，まずは説明変数が2つある重回帰分析を考えてみましょう。予測値$\hat{y}$は次のように表現できます。
\[ \hat{y}= \beta_0 + \beta_1x_1 + \beta_2x_2 \]
これを空間のイメージで考えると，右辺は説明変数$\bm{x} =(\bm{1}, \bm{x}_1,\bm{x}_2) \in \mathbb{R}^3$に係数ベクトル$\bm{\beta} = (\beta_0,\beta_1,\beta_2)' \in \mathbb{R}^3$をかけてたことによる空間が広がっていますね。$\bm{X}$が張る空間を$\bm{W(X)}$とすると，次のように表すことができます。
\[ \bm{W(X)} =\{ \beta_0 \bm{1} + \beta_1\bm{x}_1 + \beta_2 \bm{x}_2 \mid \bm{\beta} \in \mathbb{R}^3\} \]

$\hat{y}$はこの空間の中にあるベクトルで，実際のデータベクトル$\bm{y}$はこの右辺が張る空間に射影されたものとみることができます(図\ref{RegressionModel})。

\begin{figure}
    \centering
    \begin{tikzpicture}[scale=0.7]
        \fill[cyan!10] (-3,-4)--(-1,4)--(7,4)--(6,-4)--cycle; %四角形を赤:青=30:70の色で塗る
        \coordinate[label=above left:$o$] (o)at (0,0);
        \coordinate (y) at (2.5,5);
        \coordinate (x1) at (4,2);
        \coordinate (x2) at (3,-3);
        \coordinate (haty) at (2.5,0);
        \coordinate (b1) at ($(o)!(haty)!(x1)$);
        \coordinate (b2) at ($(o)!(haty)!(x2)$);
        \draw[->, >=Latex,very thick] (o) -- (y);
        \draw[->, >=Latex,thick] (o) -- (x1);
        \draw[->, >=Latex,thick] (o) -- (x2);
        \draw[->, >=Latex,very thick] (o) -- (haty);
        \draw[dashed] (y) -- node[midway,right]{$\bm{e}$}(haty);
        \draw[dashed] (haty) --(b1);
        \draw[dashed] (haty) --(b2);
        \node[above] at (y) {$\bm{y}$};
        \node[below] at (haty) {$\bm{\hat{y}}$};
        \node[right] at (x1) {$\bm{x1}$};
        \node[right] at (x2) {$\bm{x2}$};
        \node[above] at (b1) {$\beta_1$};
        \node[below left] at (b2) {$\beta_2$};
        \node[below left] at (7,4) {$\bm{W(X)}$};
    \end{tikzpicture}
    \caption{説明変数の空間に被説明変数ベクトルを射影する}
    \label{RegressionModel}
\end{figure}

セクション\ref{ProjectionMatrix}(Pp.\pageref{ProjectionMatrix})でみたように，射影行列$\bm{P}$は基底ベクトルを使って$\bm{P} = \bm{A}(\bm{A}'\bm{A})^{-1}\bm{A}'$と表されるということでしたから，回帰分析の場合は
\[ \bm{\hat{y}} = \bm{X}(\bm{X}'\bm{X})^{-1}\bm{X}' \bm{y} \]
として求められることになります。


ちなみに誤差$\bm{e}=\bm{y} - \bm{\hat{y}}$は，被説明変数$\bm{y}$から空間$\bm{W(X)}$に射影した時の垂線の長さです。推薦ですから，説明変数が張る空間とは直角に交わっています。すなわち，\underline{誤差は説明変数との相関がゼロ}である，という回帰分析の特徴も一気に理解することができます。

また，垂線ですので，被説明変数ベクトルと説明変数が張る空間との距離が最短になっていることもわかります。
平面$\bm{W(X)}$上の別のベクトル$\bm{P}$を考えて$\bm{P}$と$\bm{y}$との距離を考えるとします。


\begin{figure}[h]
    \centering
    \begin{tikzpicture}[scale=1]
        \coordinate (o) at (0,1);
        \coordinate (E) at (10,2);
        \coordinate (y) at (4,4);
        \coordinate (haty) at ($(o)!(y)!(E)$);
        \coordinate (P) at ($(o)!0.75!(E)$);
        \draw[cyan!10,line width=2pt] (o) -- (E);
        \draw[dashed] (y) --node[midway,left]{$\bm{e}$} (haty);
        \draw[dashed] (y) --node[midway,right]{$\bm{y}-\bm{P}$} (P);
        \node[above] at (y) {$\bm{y}$};
        \node[below] at (haty) {$\bm{\hat{y}}$};
        \node[below] at (P) {$\bm{P}$};
        \node[auto] at (E) {$\bm{W(X)}$};
    \end{tikzpicture}
    \caption{垂線が最短距離であることを考える}
    \label{minQ}
\end{figure}

\[
    \begin{array}{lll}
        \| \bm{y} - \bm{P} \|^2 & = & \| \bm{y}-\bm{\hat{y}} + \bm{P} - \bm{\hat{y}} \|^2                                       \\
                                & = & \|\bm{y} - \bm{\hat{y}}\|^2 + \|\bm{P}-\bm{\hat{y}}\|^2 + 2\bm{e}'(\bm{P} - \bm{\hat{y}}) \\
                                & = & \| \bm{e} \|^2 + \|\bm{P}-\bm{\hat{y}}\|^2
    \end{array}
\]
つまり，別のベクトル$\bm{P}$はどうやっても$\bm{e} = \bm{y}-\bm{\hat{y}}$以上になるので，$\bm{e}$が最短だということになります。この$\bm{e}$のノルムを最小にするのが\keytermL{最小二乗法}{さいしょうにじょうほう}ですが，距離を最小にするので二乗なんだな，と考えることができます。

具体的な係数ベクトル$\bm{\beta}$は，空間$\bm{W(X)}$と残差$\bm{e} = \bm{y} - \bm{\hat{y}}$が直交することから，
\[ \bm{X}'\bm{e} = 0　\]
\[ \bm{X}'(\bm{y} - \bm{X\beta}) = 0 \]
\[ \bm{X}'\bm{y} - \bm{X}'\bm{X}\bm{\beta} = 0 \]
\[ \bm{X}'\bm{X}\bm{\beta} = \bm{X}'\bm{y} \]
となりますから，$\bm{\beta}$だけにするには$\bm{X}'\bm{X}$に逆行列があればよく，
\[ \bm{\beta} = (\bm{X}'\bm{X})^{-1}\bm{X}'\bm{y} \]
ということになります。そもそも$\bm{\hat{y}} = \bm{X\beta}$だったわけですから，射影行列との関係$\bm{\hat{y}} = \bm{X}(\bm{X}'\bm{X})^{-1}\bm{X}' \bm{y} $とも整合的であることがわかりますね。
\subsection{多重共線性の問題}
重回帰分析において\keyterm{多重共線性}{たじゅうきょうせんせい}{multicolinearity}に気をつけなければならない，というのは聞いたことがあると思います。

多重共線性は，説明変数同士の相関が高いことが原因でおこることで，なんか良くないんだ，というイメージはお持ちかと思います。さて何が悪いのでしょう?
これを考えるときに，空間的なイメージからアプローチしてみましょう。説明変数同士の相関が高いこと，というのがありましたが，相関が高いことを空間のイメージでいうと，ベクトルの$\cos \theta$が0ではないということです。相関が最高潮に高い場合は，ベクトルがぴたりと重なってしまいますから，これは一次独立でなくなり，空間が潰れてしまうことになります。

今回の例では2つの説明変数が作る空間があって，そこに被説明変数ベクトルの影を落としたのでした。図\ref{RegressionModel}にもあるように，説明変数ベクトルの座標が偏回帰係数になります。このとき，2つの説明変数が実は一次従属関係にあったとすると，2本のベクトルではなく1本のベクトル空間になります。そんな状況で2つの回帰係数を求めることはできないですね。

重回帰分析が実践的に使われるシーンでは，2つと言わずもっと多くの説明変数が使われることがあります。$n$人分のデータ，$p$個の説明変数があったとして，行列$\bm{X}$は$n\times (p+1)$になりますが，説明変数のなかに一次従属な変数があるとこの行列$\bm{X}$のランクが$rank(\bm{X}) < (p+1)$になっていることになります\footnote{行列のランクは行数，列数の小さい方であり，一般にデータ数$n$のほうが変数の数$p$よりも大きい($p\ll n$)ので，ここでのランクは$p+1$で考えます。ちなみに$\ll$は「十分にちいさい」という意味です。}。行列の積とランクの関係から，
\[ rank(\bm{X}'\bm{X}) \le rank(\bm{X}) \]
ですが，$p+1 \times p+1$の正方行列$\bm{X}'\bm{X}$が逆行列を持つ条件は$rank(\bm{X}'\bm{X}) = p+1$であることで，ランク落ち$rank(\bm{X}'\bm{X}) < p+1$が生じていると逆行列が持てません。行列$(\bm{X}'\bm{X})$は射影するときに逆行列にして使いますので，ここで計算上のエラーが生じます。

逆行列の計算は，行列のサイズが大きくなっていくと非常に計算が大変になって，近似的に解いていくことになります。解析的には答えがないはずの行列に対して，無理矢理近似的な数値を算出させたところで，その数字が信用できるはずもないことはおわかりいただけると思います。もちろん実際のデータで，変数同士が一次従属になることは稀です\footnote{変数を生成して説明変数に加えると，こうした問題が生じがちです。例えば身長の単位をcmからmに変えて両方とも投入してしまうといった人為的なミスというケースです。}が，尺度の合計得点と因子得点といった，相関がごく高い変数の両方を使用変数にしてしまうと，ランク落ちの様相を呈してきます。もちろん偽相関など見せかけ的に高い関係があってもいけません。

\begin{Untiku}{Rでやってみて考える}{}
    実際に数値例をつかってみてみればいいのですが，逆行列の計算は面倒なので，Rをつかって計算させてみたいと思います。
    Rで逆行列を求めるには，\verb|solve|という関数があるので，これを使います。
    \begin{lstlisting}[language=C++,label={Multico1},caption={フルランクの行列}]
    A <- matrix(c(1,3,5,2,4,3,3,1,2,2,1,6,3,5,4),ncol=3)
    M <- t(A) %*% A
    solve(M)
    \end{lstlisting}
    まず適当な行列$\bm{A}$を作り，$\bm{M} = \bm{A}'\bm{A}$を構成します。今回の例では，
    \[ \bm{A} = \begin{pmatrix}1 & 3 & 1\\ 3&3&6\\5&1&3\\2&2&5\\4&2&2\end{pmatrix},\bm{M} = \begin{pmatrix}55&29&60\\29&27&42\\60&42&87\end{pmatrix}\]
    であり，逆行列は
    \[ \bm{M}^{-1}=\begin{pmatrix} 0.0734 & -0.0003 & -0.0504\\-0.0003 & 0.1487 & -0.0715\\-0.0504&-0.0715&0.0808\end{pmatrix}\]
    となりました。

    ここでわざとランク落ちを生じさせます。$\bm{A}$の3列目を，$\bm{A}$の2列目で上書きしてしまうのです。
    \begin{lstlisting}[language=C++,label={Multico2},caption={ランクおちの行列}]
    A[,3] <- A[,2]
    M <- t(A) %*% A
    solve(M)
    \end{lstlisting}
    こうすると，Rはエラーを返します。
    \begin{verbatim}
    solve.default(M) でエラー: 
    Lapack routine dgesv: システムは正確に特異です: U[3,3] = 0 
    \end{verbatim}
    だそうです。

    さて，ではランク落ちかけ，という程度ならどうでしょう。$\bm{A}$の3列目を$\bm{A}$の2列目で上書きするときに，少し数字をいじってしまいます。
    \begin{lstlisting}[language=C++,label={Multico2},caption={ランクおちかけ}]
    A[,3] <- A[,2]+rnorm(5)*0.001
    M <- t(A) %*% A
    solve(M)
    \end{lstlisting}
    今回は微小な誤差を追加しましたので，行列は
    \[
        \bm{A} =
        \begin{pmatrix}
            1 & 3 & 3.0005855 \\
            3 & 3 & 3.0007095 \\
            5 & 1 & 0.9998907 \\
            2 & 2 & 1.9995465 \\
            4 & 2 & 2.0006059 \\
        \end{pmatrix},
        \bm{M} =
        \begin{pmatrix}
            55.00000 & 29.00000 & 29.00368 \\
            29.00000 & 27.00000 & 27.00408 \\
            29.00368 & 27.00408 & 27.00816 \\
        \end{pmatrix}
    \]
    です。ほとんど変わらない大きさですが，これの逆行列は，
    \[
        \bm{M}^{-1} =
        \begin{pmatrix}
            0.0430067   & -36.95954 & 36.90777 \\
            -36.9595434 & 1260304   & -1260074 \\
            36.9077732  & -1260074  & 1259844  \\
        \end{pmatrix}
    \]
    という無茶苦茶な大きさになってしまいました。
\end{Untiku}

多重共線性が問題になるシーンは，計算がうまくいっていないかも，という数値計算の話だけではありません。データ解析の多くのシーンでは，求めたい回帰係数は母集団のそれであり，毎回のサンプルごとどれぐらい変わるかということを，\keytermL{標準誤差}{ひょうじゅんごさ}で表現します。多重共線性が生じて計算上の不確実さがあると，これが標準誤差の大きさに跳ね返ってきます。

図\ref{Multico}には，被説明変数が$\bm{y}$と$\bm{y}'$の2本描かれています。違うサンプルを取ったことで少し違う値が得られたというイメージです\footnote{この素晴らしい幾何学的な説明のアイデアは，\citet{Haebara200206}によるものです。}。説明変数ベクトルが張る空間に落とした影から，説明変数ベクトルに垂線を引いたところが係数になります。右図にあるように，説明変数同士に相関があると，ちょっと影がずれるだけで係数の符号まで変わってしまいます。このような不安定さが簡単に生じるので，多重共線性があった場合は係数が信用ならないわけです。


\begin{figure}
    \centering
    \begin{tikzpicture}[scale=0.65]
        \coordinate (o) at (0,0);
        \coordinate (b1a) at (-5,5);
        \coordinate (b1b) at (5,-5);
        \coordinate (b2a) at (-5,-5);
        \coordinate (b2b) at (5,5);
        \coordinate (y1) at (0.3,4);
        \coordinate (y2) at (-0.3,4);
        \coordinate (yhat1) at (0.5,-3);
        \coordinate (yhat2) at (-0.5,-3);
        \coordinate (y1b1) at ($(b1a)!(yhat1)!(b1b)$);
        \coordinate (y1b2) at ($(b2a)!(yhat1)!(b2b)$);
        \coordinate (y2b1) at ($(b1a)!(yhat2)!(b1b)$);
        \coordinate (y2b2) at ($(b2a)!(yhat2)!(b2b)$);
        \draw[->,>=Latex,thick] (b1a)--(b1b);
        \draw[->,>=Latex,thick] (b2a)--(b2b);
        \draw[->,>=Latex,very thick] (o) -- (y1);
        \draw[->,>=Latex,very thick] (o) -- (y2);
        \draw[->,>=Latex,very thick,gray!80] (o) -- (yhat1);
        \draw[->,>=Latex,very thick,gray!80] (o) -- (yhat2);
        % \draw[dashed] (y1) -- (yhat1);
        % \draw[dashed] (y2) -- (yhat2);
        \draw[dashed] (yhat1) -- (y1b1);
        \draw[dashed] (yhat1) -- (y1b2);
        \draw[dashed] (yhat2) -- (y2b1);
        \draw[dashed] (yhat2) -- (y2b2);
        \node[below right] at (b1b) {$\bm{b}_1$};
        \node[below right] at (b2b) {$\bm{b}_2$};
        \node[above right] at (y1) {$\bm{y}$};
        \node[above left] at (y2) {$\bm{y}'$};
        \node[above right] at (y1b1) {$\beta_1$};
        \node[above left] at (y1b2) {$\beta_2$};
        \node[above right] at (y2b1) {$\beta_1'$};
        \node[above left] at (y2b2) {$\beta_2'$};
    \end{tikzpicture}
    \begin{tikzpicture}[scale=0.65]
        \coordinate (o) at (0,0);
        \coordinate (b1a) at (-5,4);
        \coordinate (b1b) at (5,-4);
        \coordinate (b2a) at (-5,0.3);
        \coordinate (b2b) at (5,-0.3);
        \coordinate (y1) at (0.3,4);
        \coordinate (y2) at (-0.3,4);
        \coordinate (yhat1) at (0.5,-3);
        \coordinate (yhat2) at (-0.5,-3);
        \coordinate (y1b1) at ($(b1a)!(yhat1)!(b1b)$);
        \coordinate (y1b2) at ($(b2a)!(yhat1)!(b2b)$);
        \coordinate (y2b1) at ($(b1a)!(yhat2)!(b1b)$);
        \coordinate (y2b2) at ($(b2a)!(yhat2)!(b2b)$);
        \draw[->,>=Latex,thick] (b1a)--(b1b);
        \draw[->,>=Latex,thick] (b2a)--(b2b);
        \draw[->,>=Latex,very thick] (o) -- (y1);
        \draw[->,>=Latex,very thick] (o) -- (y2);
        \draw[->,>=Latex,very thick,gray!80] (o) -- (yhat1);
        \draw[->,>=Latex,very thick,gray!80] (o) -- (yhat2);
        % \draw[dashed] (y1) -- (yhat1);
        % \draw[dashed] (y2) -- (yhat2);
        \draw[dashed] (yhat1) -- (y1b1);
        \draw[dashed] (yhat1) -- (y1b2);
        \draw[dashed] (yhat2) -- (y2b1);
        \draw[dashed] (yhat2) -- (y2b2);
        \node[below right] at (b1b) {$\bm{b}_1$};
        \node[below right] at (b2b) {$\bm{b}_2$};
        \node[above right] at (y1) {$\bm{y}$};
        \node[above left] at (y2) {$\bm{y}'$};
        \node[above right] at (y1b1) {$\beta_1$};
        \node[above right] at (y1b2) {$\beta_2$};
        \node[above right] at (y2b1) {$\beta_1'$};
        \node[below left] at (y2b2) {$\beta_2'$};
    \end{tikzpicture}
    \caption{左;サンプルが変わると少し係数がぶれる。右；サンプルが変わっただけで$\beta_2$の符号も変わってしまう。}
    \label{Multico}
\end{figure}

\subsection{実験計画による空間の分解}

\section{分散共分散行列}



\subsection{回転行列と分散}



\end{document}